\documentclass[11pt]{article}
\usepackage[utf8]{inputenc}
\usepackage{amsmath,amssymb,amsthm}
\usepackage[margin=1in]{geometry}

\newtheorem{lemma}{Lemma}
\newtheorem{theorem}{Theorem}

\title{Problem 506: Clock Sequence}
\author{}
\date{}

\begin{document}
\maketitle

\section{Problem Statement}

Consider the infinite periodic digit stream
\[
123432123432123432\cdots
\]
and read it from left to right. Define $v_n$ to be the shortest consecutive block of unread digits whose digit sum is exactly $n$. The beginning of the sequence is
\[
1,\ 2,\ 3,\ 4,\ 32,\ 123,\ 43,\ 2123,\ 432,\ 1234,\ 32123,\ldots
\]
Let
\[
S(N)=\sum_{n=1}^{N} v_n.
\]
It is given that $S(11)=36120$. Find $S(10^{14}) \bmod 123454321$.

\section{Mathematical Analysis}

Let $c=123432$ be the fundamental 6-digit cycle. Its digit sum is $15$.

\begin{lemma}
The first fifteen terms are
\[
\begin{aligned}
&v_1=1,\quad v_2=2,\quad v_3=3,\quad v_4=4,\quad v_5=32,\\
&v_6=123,\quad v_7=43,\quad v_8=2123,\quad v_9=432,\quad v_{10}=1234,\\
&v_{11}=32123,\quad v_{12}=43212,\quad v_{13}=34321,\quad v_{14}=23432,\quad v_{15}=123432.
\end{aligned}
\]
\end{lemma}

\begin{proof}
Since all digits are positive, for each target sum $n$ there is a unique shortest unread block whose digit sum is $n$. A direct simulation yields the displayed values.
\end{proof}

\begin{lemma}
After the first fifteen terms have been read, the stream is again aligned with the start of the cycle $c$.
\end{lemma}

\begin{proof}
The total number of digits consumed by the fifteen terms is
\[
1+1+1+1+2+3+2+4+3+4+5+5+5+5+6 = 48 = 8 \cdot 6,
\]
which is a multiple of the cycle length. Hence the next unread digit is again the first digit of $c$.
\end{proof}

\begin{theorem}
The pattern of starting phases repeats with period $15$. Hence for every $q \ge 0$ and every $1 \le r \le 15$, the block $v_{15q+r}$ starts in the same phase as $v_r$.
\end{theorem}

\begin{proof}
By the previous lemma the stream is realigned after one block of fifteen terms. Repeating the same argument inductively proves the claim for every subsequent 15-term block.
\end{proof}

For each residue $r \in \{1,\dots,15\}$ let $p_r$ be the starting phase of $v_r$, let $C_r$ be the 6-digit rotation of $c$ starting at phase $p_r$, and let $U_r$ be the terminal partial block needed to realize the residue $r \bmod 15$.

Writing
\[
n=15q+r \qquad (1 \le r \le 15),
\]
we obtain
\[
v_n = \underbrace{C_r C_r \cdots C_r}_{q\text{ copies}}\, U_r.
\]
Therefore, modulo $M=123454321$,
\[
v_n \equiv A_r G(q) + B_r \pmod M,
\]
where
\[
A_r \equiv C_r 10^{L_r} \pmod M,\qquad B_r \equiv U_r \pmod M,
\]
and
\[
G(q)=1+10^6+10^{12}+\cdots+10^{6(q-1)}.
\]

\begin{theorem}
The contribution of all complete 15-term blocks reduces to geometric sums.
\end{theorem}

\begin{proof}
The coefficients $A_r,B_r$ depend only on the residue class $r$, not on the block index. Hence
\[
\sum_{q=0}^{Q-1} v_{15q+r}
\equiv
A_r \sum_{q=0}^{Q-1} G(q) + Q B_r \pmod M.
\]
Since $G(q)$ is itself a geometric sum in the ratio $10^6$, the outer sum is again expressible by geometric-series identities.
\end{proof}

\section{Algorithm}

\begin{enumerate}
\item Simulate the first fifteen terms once and record the phase data $(p_r,U_r,L_r)$.
\item Form the rotations $C_r$ and coefficients $A_r,B_r$.
\item Let $N=10^{14}$, $Q=\lfloor N/15 \rfloor$, and $R=N \bmod 15$.
\item Sum the $Q$ complete blocks using the geometric-series formula above.
\item Add the first $R$ terms of the next block directly.
\end{enumerate}

\begin{theorem}
The algorithm computes $S(10^{14}) \bmod 123454321$.
\end{theorem}

\begin{proof}
Lemma 1 gives the exact first-block data. Lemma 2 and the first theorem imply 15-periodicity of the starting phases. The decomposition $v_{15q+r}=C_r^q U_r$ is exact because each full cycle contributes digit sum $15$ and returns to the same phase. The geometric-sum reduction therefore computes the contribution of every complete block exactly, and the remaining $R$ terms are added explicitly. Hence every term $v_n$ with $1 \le n \le 10^{14}$ is counted once and only once.
\end{proof}

\section{Complexity}

The precomputation is constant-size. The final evaluation uses a constant number of modular exponentiations and geometric-sum evaluations.

\begin{itemize}
\item Time: $O(\log N)$.
\item Space: $O(1)$.
\end{itemize}

\section{Answer}

\[
\boxed{18934502}
\]

\end{document}
